\documentclass[10pt,a4paper]{amsart}
\usepackage[margin=19mm]{geometry}
\usepackage{amsmath,amssymb,mathtools}
\usepackage[scr=rsfs,cal=euler]{mathalfa}
\usepackage{xfrac}
\usepackage[unicode,hidelinks,bookmarks=false]{hyperref}
\newcommand{\ie}{i.e.\ }
\newcommand{\cf}{cf.\ }
\newcommand{\resp}{respectively\ }
\newcommand{\Subsection}{\subsection}
\providecommand{\nobreakdash}{\nobreak\hbox{-}}
\setlength{\emergencystretch}{2em}
\multlinegap0pt
\usepackage{bm}
\usepackage{bbm}
\usepackage{dsfont}
\usepackage{xspace}
\usepackage{cancel}
\usepackage{mathtools}
\usepackage{amsthm}
\makeatletter
\@ifundefined{theorem}{\newtheorem{theorem}{Theorem}}{}
\@ifundefined{lemma}{\newtheorem{lemma}[theorem]{Lemma}}{}
\@ifundefined{prop}{\newtheorem{prop}[theorem]{Proposition}}{}
\makeatother
\newcommand{\Cc}{\mathcal{C}}
\newcommand{\Dc}{\mathcal{D}}
\newcommand{\E}{\mathcal{E}}
\newcommand{\Gal}{\mathrm{Gal}}
\newcommand{\N}{\mathbb{N}}
\newcommand{\Q}{\mathbb{Q}}
\newcommand{\Z}{\mathbb{Z}}
\newcommand{\im}{\mathrm{im}}
\title[The Smallest Invariant Factor of Elliptic Curves, and Coinci]{The Smallest Invariant Factor of Elliptic Curves, and Coincidences}
\author{Alexander Milner, Jack Shotton}
\date{Source: arXiv:2604.21601v2 (CC BY 4.0)}
\begin{document}
\maketitle

\noindent\emph{Source and licence.} The Smallest Invariant Factor of Elliptic Curves, and Coincidences. Version arXiv:2604.21601v2 (CC BY 4.0), distributed under the Creative Commons Attribution 4.0 International licence, as recorded in the arXiv licence metadata for this exact version and shown on the abstract page https://arxiv.org/abs/2604.21601v2. Full rights notice: licenses/arxiv-2604.21601-CC-BY-4.0.txt.\par\smallskip

\noindent\textit{Editorial scope.} Counted unit: the Theorem whose source label is \texttt{T4b} in the article (source0.tex lines 468--472, source label \texttt{T4b}), delivered together with the article's own complete original proof of it (source0.tex lines 473--523, one \texttt{proof} environment of 51 lines). No mathematics is written, repaired or re-derived: statement and proof are the article's own text line for line. Required context retained uncounted: prop:Cejiff (lines 327--329); lem:cyclotomic-lemma (lines 456--461). Every definition, display and label of those lines is retained unchanged. Omitted from this package and not redistributed: 1--326, 330--455, 462--467, 524--917. Nothing inside a retained excerpt is omitted or altered: every retained body is delivered exactly as the article's lines are written, so no omission receipt of any kind accompanies this package. Citations: the retained excerpts carry no citation command at all, so no bibliography entry of the article is transcribed and no reference is redistributed. Reference adaptation: none was needed and none was made; every reference inside the delivered unit resolves inside it, so no unresolved reference or citation is produced. Printed slips kept literal: no printed word, symbol, hypothesis or number of the counted statement or of its proof was altered. Layout and class: the article's own class is \texttt{amsart} with packages including amsmath, amsthm, mathtools, tikz-cd, xurl, fontenc, tabularx, colortbl, enumitem, amssymb, stmaryrd, xcolor, verbatim, float; the wrapper is the lane's pinned \texttt{amsart} editorial wrapper at 10pt on A4, which restates the theorem-like environments the retained text needs and adds the article's own macro definitions that the retained text needs (\texttt{\textbackslash Cc}, \texttt{\textbackslash Dc}, \texttt{\textbackslash E}, \texttt{\textbackslash Gal}, \texttt{\textbackslash N}, \texttt{\textbackslash Q}, \texttt{\textbackslash Z}, \texttt{\textbackslash im}), with extra wrapper packages (bm, bbm, dsfont, xspace, cancel, mathtools). The delivered numbering is the unit's own. Assets: no retained span contains a figure environment, a graphics-inclusion command, a PDF-page inclusion, a table or a TikZ picture; no image file is copied into this package, the package declares no asset dependency and no image file is redistributed with it. Licence: version arXiv:2604.21601v2 of this article is distributed under the Creative Commons Attribution 4.0 International licence, as recorded in the arXiv licence metadata for this exact version and shown on the abstract page https://arxiv.org/abs/2604.21601v2. A full rights notice is delivered at licenses/arxiv-2604.21601-CC-BY-4.0.txt. Characters: every retained excerpt is pure ASCII, so the delivered engine, XeLaTeX with its default Latin Modern text font, reports no missing character.
\begin{prop}\label{prop:Cejiff} Suppose that $j \mid m$ and let $S = \{p : p \mid m/j\}$. Then $\Cc_{\E, j} = 0$ if and only if 
\[H \cap \Gamma_j \bmod m = \bigcup_{p \in S}H \cap \Gamma_{jp} \bmod m.\]
\end{prop}

\begin{lemma}\label{lem:cyclotomic-lemma}
    Suppose that $j, k \in \N$ with $jk \mid m_\E$. Then 
    \[\Q(\E[j]) \cap \Q(\zeta_{jk}) = \Q(\zeta_j)\]
    if and only if 
    \[\Dc_{jk}(H \cap \Gamma_j) = (1 + j\Z/jk\Z) \subset (\Z/jk\Z)^\times.\]
\end{lemma}

\begin{theorem} \label{T4b}
    Let $j \mid m_\E$ be a natural number, let $R$ be the product of $j$ with the odd prime factors of $m_\E/j$.
    
    If $\mathbb{Q}(\E[j]) \cap \mathbb{Q}(\zeta_{R}) =\mathbb{Q}(\zeta_{j}) $ and $\mathbb{Q}(\E[2j]) \neq \mathbb{Q}(\E[j])$, then $\mathcal{C}_{\E,j}>0$.
\end{theorem}

\begin{proof}
    Write $m_\E/j = \prod_{i=1}^r p_i^{a_i}$ with $a_i \in \N$ and $p_i$ distinct primes. Let $m_i = v_{p_i}(m_\E)$ and $j_i = v_{p_i}(j)$, so that $a_i = m_i - j_i$. 

    For any integer $n\ge 0$, set 
    \[K_n(p_i) = \begin{cases}
        I + p_i^nM_2(\Z/p_i^{m_i}\Z) & \text{if $n > 0$, and}\\
        GL_2(\Z/p_i^{m_i}\Z) & \text{if $n = 0$.}
    \end{cases}\]
    Similarly, set 
     \[U_n(p_i) = \begin{cases}
        1 + p_i^n\Z/p_i^{m_i}\Z & \text{if $n > 0$, and}\\
        (\Z/p_i^{m_i}\Z)^\times & \text{if $n = 0$.}
    \end{cases}\]

    Then $H \cap \Gamma_j \subset \prod_{i=1}^r K_{j_i}(p_i)$ and to show that $\Cc_{\E,j} > 0$ it is enough, by Proposition~\ref{prop:Cejiff}, to show that 
    \[H \cap \Gamma_j \ne \bigcup_{i=1}^r H \cap \Gamma_{p_ij}.\]
    This is equivalent to showing that there is $h \in H \cap \Gamma_j$ such that $\pi_i(h) \ne e$ for all $1 \le i \le r$, where 
    \[\pi_i : H \cap \Gamma_j \to K_{j_i}(p_i)/K_{j_i + 1}(p_i)\]
    is the projection.

    Suppose first that $2 \not \in \{p_1, \ldots, p_r\}$. By our hypothesis and Lemma~\ref{lem:cyclotomic-lemma}, 
    \[\det : H \cap \Gamma_j \to \prod_{i=1}^r U_{j_i}(p_i)/U_{j_i+1}(p_i)\]
    is surjective. Since each factor in this product is nontrivial, there is $h \in H \cap \Gamma_j$ whose image in every factor is not the identity. This is the element $h$ we need.

    Otherwise, suppose without loss of generality that $p_1 = 2$. The image $G_1$ of $H\cap \Gamma_j$ in $K_{j_1}(2)/K_{j_1 + 1}(2)$ is isomorphic to $\Gal(\Q(\E[2j])/\Q(\E[j]))$, and hence nontrivial by hypothesis.

    Let $G_2 = \prod_{i=2}^rU_{j_i}(p_i)/U_{j_i + 1}(p_i)$ and consider the homomorphism 
    \[\theta : H \cap \Gamma_j \to G_1 \times \left(\prod_{i=2}^r U_{j_i}(p_i)/U_{j_i + 1}(p_i)\right) = G_1 \times G_2\]
    and let $T = \im(\theta)$. It is enough to show that there is $(t_1, \ldots, t_r) \in T$ such that no $t_i$ is the identity.
    
    The projection $p_1$ of $T$ onto $G_1$ is surjective by definition, while the projection $p_2$ of $T$ onto $G_2$ is
    surjective by Lemma~\ref{lem:cyclotomic-lemma}. Let $N_1 = \ker(p_2) \subset G_1$ and let
    $N_2 = \ker(p_1) \subset G_2$. By Goursat's Lemma, $N_1 \times N_2 \subset T$ and the image of $T$ in
    $G_1/N_1 \times G_2 / N_2$ is the graph of an isomorphism $\phi : G_2/N_2 \to G_1/N_1$.

    Suppose that $x = (x_2, \ldots, x_r) \in G_2$ with $x_i \ne e$ for every $2 \le i \le r$. Then $\phi(x) = e$
    otherwise $(\phi(x), x_2, \ldots, x_r)$ would be the required element $t$. So $N_2$ contains every element of this
    form.

    Now let $x = (x_2, \ldots, x_r) \in G_2$ be arbitrary. For each $2 \le i \le r$ we may choose $y_i$ and $z_i$ in
    $U_{j_i}(p_i)/U_{j_i+1}(p_i)$ such that $y_i z_i = x_i$ and $y_i, z_i \ne e$, \emph{unless} $p_i = 3$, $j_i = 0$ and
    $x_i = 2$. Suppose for now that this does not happen. Then taking $y = (y_2, \ldots, y_r)$ and
    $z = (z_2, \ldots, z_r)$, we have $x = yz$ and, by the previous paragraph, $y, z \in N_2$. So $x \in N_2$. But then,
    as $x$ was arbitrary, $N_2 = G_2$ and so $T = G_1 \times G_2$, in which case we can easily choose $t \in T$ that is
    not the identity in any factor.

    If there is some $i$ with $p_i = 3$, $j_i = 0$ and $x_i = 2$ then, without loss of generality, $i = 2$. We then take
    $w = (2,w_2, \ldots, w_r)$, $y = (2, y_2, \ldots, y_r)$ and $z = (2, z_2, \ldots, z_r)$ where all
    $w_i, y_i, z_i \ne e$ and $w_iy_iz_i = x_i$ (for example, take $w_i \ne e$ arbitrary and then find $y_i, z_i \ne e$
    such that $y_i z_i = x_i w_i^{-1}$). Then $w, y, z \in N_2$ and $x = wyz$; we conclude as before.
\end{proof}

\end{document}
